\clearpage
\fancyhead[L]{\small TAVOLE DELLE PARALLASSI.}
\hypertarget{native-604-lower}{}
\begin{center}\small Göttingische gelehrte Anzeigen. No. 108. Pages 1073..1077. 7 July 1817.\\[3pt]\rule{.72\linewidth}{.3pt}\end{center}

\emph{Tavole delle parallassi di altezza, di longitudine e di latitudine calcolate dagli astronomi dell' osservatorio dell' università Gregoriana nel collegio romano.} Rome, 1816. Nella stamperia de Romanis. XXXIV and 122 pages in folio.

\begingroup\brokenpenalty=10000\relax
Occultations of fixed stars by the Moon and solar eclipses are of such great importance to astronomy and geography that any contribution toward easing the calculations relating to them deserves to be gratefully received. The so-called method of the nonagesimal is still the one most used by astronomers to calculate those phenomena, and calculating the Moon's parallax in longitude and latitude forms a principal part of that method. The present tables are intended to shorten this calculation of the parallaxes. The plan originated with \textsc{Andreas Conti}, one of the astronomers at the Roman Observatory. In carrying out this extensive work he was assisted by his colleagues \textsc{Joseph Calandrelli} and \textsc{Jacob Reichenbach}.

To enable our readers to form an idea of how these tables are arranged and used and of the gain to be achieved thereby, we must examine each of them more closely. There are actually \emph{three} tables. The first contains 93 pages, and the second and third 12 pages each. A separate sheet also contains a few small auxiliary tables. The first table contains only the products of half the horizontal parallax and the cosines of all angles. It thus has two arguments. One argument is the horizontal parallax, whose 50 values run from $53'\,20''$ to $61'\,30''$ in steps of 10 seconds. The second argument consists of all angles of the quadrant in steps of 10 minutes. Each page is divided in two. The larger half contains the products just mentioned, while the smaller half gives, under the title Supplement to the First Table, the products of the same cosines and half the horizontal parallaxes of the planets (from $1''$ to $40''$). It further gives the products of these cosines and each individual minute from $1'$ to $9'$, and each multiple of ten minutes from $10'$ to $90'$. As one can see, then, space has not been spared.

\SourceLeaf{605}The second table likewise has two arguments. It gives us the value of the product $\alpha(\sec L - 1)$ for the arguments $L$ and $\alpha$, the former taken from 0 to $6^{\circ}$ at every multiple of ten minutes, the latter from 0 to $60'\,30''$ in half-minute steps and then also from $1''$ to $35''$ at every odd second.

The third table contains 12 fixed stars that are frequently occulted by the Moon, the value of an auxiliary angle $\varphi$ that depends partly on the latitude of the fixed star and partly on an angle $D - \tfrac{1}{2}\Pi$ to be explained below. The values of this angle run through the whole quadrant from 20 to 2c minutes.

The three small auxiliary tables also contain the difference between the tangent and the arc, the conversion of seconds into decimal parts of a minute, and the value of the product, $\operatorname{tang} \tfrac{1}{2}(D+\tfrac{1}{2}\Pi)\cdot\sin L'$, where the first argument $D+\tfrac{1}{2}\Pi$ comprises all angles of the quadrant and the second, $L'$, the angles from 0 to $6^{\circ}$ in one-degree steps.

This is the content of the tables. It remains for us to explain how they are to be used.

The parallax in altitude is obtained at once by doubling the result from the first table when one enters it with the horizontal parallax and the apparent altitude. If the true altitude is given instead, one will obtain the desired quantity by successive approximation. The same remark also applies to what follows insofar as the arguments to be used are not given directly.

The parallax in longitude $\Pi$ is determined, always with sufficient accuracy, by the formula
\[
\Pi = \dfrac{P \sin h \sin(D+\Pi)}{\cos L}
\]
where $P$ denotes the horizontal parallax, $h$ the altitude of the nonagesimal, $D$ the true longitude less the longitude of the nonagesimal, and $L$ the true latitude. If $\alpha$ is the difference between the results from the first table when one enters it first with $P$ and $D-h+\Pi$, and then with $P$ and $D+h+\Pi$, then $\Pi$ will equal the sum of $\alpha$ and the result from the second table when one enters it with the arguments $\alpha$ and $L$.

The parallax in latitude $\pi$ follows with sufficient accuracy from the formula
\[
\pi = -P[\cos h \cos(L+\pi) - \sin h \sin(L+\pi)\cos(D+\tfrac{1}{2}\Pi)]
\]
or from this one
\[
\pi = -\dfrac{P\cos(h+\varphi)\cos(L+\pi)}{\cos\varphi}
\]
where $\varphi$ denotes an auxiliary angle determined by the formula
\[
\operatorname{tang} \varphi = \operatorname{tang}(L+\pi)\cdot\cos(D+\tfrac{1}{2}\Pi)
\]
Since the last formula for $\pi$ has the same form as the formula for $\Pi$, it is easy to see how the first and second tables can be used to calculate it.

The third table gives the aids $\varphi$ when the occultation of one of its twelve stars is to be calculated by \textsc{Carlini}’s method. In general, however, one finds it by combining the supplement to the first table with the first auxiliary table. Since the auxiliary angle $\varphi$ depends on the apparent latitude, the following procedure is prescribed when the true latitude is given instead. If we set, as above,
\[
P \sin h \sin(D+\Pi) = \alpha
\]
and at the same time
\SourceLeaf{606}\[
\sin(L+\pi)\,\operatorname{tang}\tfrac{1}{2}(D+\tfrac{1}{2}\Pi) = M
\]
then, with sufficient accuracy,
\[
\pi = -P\cos(h+\lambda+\pi) - \alpha M
\]
where $M$ is found from the last auxiliary table, and $\alpha$ and $P\cos(h+\lambda+\pi)$ with the help of the first table. Successive approximation is unavoidable here too.

If we are now to state our opinion of these tables frankly, we must admit that we can find no simplification of the parallax calculations in them. We consider calculation with the ordinary trigonometric tables at least as convenient as the use of these special tables. The printing and paper, incidentally, are fine, just as one would wish them to be in every numerical work.

The dedication of this work to the Pope is also noteworthy. One is glad of the support he gives the Roman Observatory. But what are we to make of the passionate zeal of the Roman astronomers against ‘the godless, who make the most shameful and dreadful misuse of astronomy, corrupt the heart, revile the true believers and mock them as senseless worshippers of the Sun’? Who understands these strange accusations, and who, without an unmistakable hint in this dedication, would have guessed that they were meant for a great geometer who had dared to put forward a hypothesis for the physical explanation of the formation of the solar system?
\par\clearpage\endgroup
